\subsection{Planteamiento del Problema}

La producción artesanal de kéfir en la Ciudad de México ha experimentado un
crecimiento notable debido a su popularidad como bebida probiótica. No obstante,
al tratarse de una actividad artesanal, el kéfir no es elaborado por una sola empresa,
sino por diversos productores independientes que aplican métodos distintos. Estos
métodos varían en aspectos clave como el tipo o la marca de leche comercial
utilizada la cual, según la Norma Oficial Mexicana NOM-155-SCFI-2012, suguiere un
pH inicial entre 6.6 y 6.8 \cite{zotero-item-119}, la cantidad de cultivos madre (búlgaros) empleados, y
las condiciones ambientales, especialmente la temperatura, que debe mantenerse
dentro de un rango recomendado de 18°C a 30°C. Esta diversidad de prácticas
genera inconsistencias tanto en los tiempos de fermentación, que pueden ir de 12
a 48 horas, como en la calidad del producto final \cite{Fitrianingsih2022}.\newline
Uno de los principales retos que enfrentan los productores es asegurar que el kéfir
sea seguro para el consumo, lo que implica identificar el momento preciso para su
recolección. Este momento ocurre cuando el pH alcanza un intervalo saludable, de
entre 4.2 y 4.6, y la temperatura se ha mantenido estable dentro del rango adecuado
durante todo el proceso. Si bien puede alcanzarse el pH requerido, si la temperatura
durante la fermentación no ha sido la apropiada, el producto puede resultar
inadecuado para el consumo, ya que las desviaciones térmicas afectan
negativamente el desarrollo de los microorganismos presentes en los búlgaros \cite{ArevaloOrtiz2007}.
De igual forma, si el kéfir se recolecta y consume antes de alcanzar el pH mínimo
(por encima de 4.6), los microorganismos no habrán completado su desarrollo,
comprometiendo sus beneficios para la salud. Por otro lado, si se consume cuando
el pH ya ha descendido por debajo del umbral recomendado (por debajo de 4.2), el
proceso fermentativo puede volverse excesivo, generando un sabor avinagrado que
indica la degradación de los microorganismos; en este estado, su consumo puede
representar un riesgo sanitario. Por lo tanto, es fundamental que el kéfir se consuma
dentro del intervalo adecuado de pH y que la temperatura permanezca dentro de
los márgenes recomendados, a fin de garantizar el desarrollo correcto de los
microorganismos y preservar sus propiedades probióticas \cite{Araujo-Osorio}.\newline
En consecuencia, el monitoreo constante del proceso fermentativo se vuelve una
tarea exigente para los productores, quienes con frecuencia recurren a métodos
empíricos para evaluar el estado del producto, como abrir el frasco para olerlo o
probarlo esporádicamente. Sin embargo, estas prácticas no solo interrumpen la
fermentación, sino que también pueden comprometer la calidad del kéfir \cite{CuevasManilla1997}.\newline
Ante esta problemática surge la siguiente pregunta: ¿Cómo desarrollar un sistema
para automatizar el monitoreo del kéfir para determinar si la fermentación alcanzo
el pH adecuado y se mantuvo en la temperatura adecuada para el consumo
humano?, ya que de mejorar este proceso de monitoreo permitirá a los productores
transitar de métodos empíricos e invasivos hacia un sistema basado en datos
estructurados y confiables.